%
% FORMATO LATEX TRABAJO FIN DE MÁSTER 
% MÁSTER PROPIO EN CIBERSEGURIDAD DE LA UGR
%
\documentclass[a4paper,11pt]{book}
%\documentclass[a4paper,twoside,11pt,titlepage]{book}
\usepackage{listings}
\usepackage[utf8]{inputenc}
\usepackage[spanish]{babel}

% \usepackage[style=list, number=none]{glossary} %
%\usepackage{titlesec}
%\usepackage{pailatino}

% Center captions
\usepackage{caption}

\decimalpoint
\usepackage{dcolumn}
\newcolumntype{.}{D{.}{\esperiod}{-1}}
\makeatletter
\addto\shorthandsspanish{\let\esperiod\es@period@code}
\makeatother


%\usepackage[chapter]{algorithm}
\RequirePackage{verbatim}
%\RequirePackage[Glenn]{fncychap}
\usepackage{fancyhdr}
\usepackage{graphicx}
\usepackage{afterpage}

\usepackage{longtable}

% Para que los links de la bibliografía no se salgan del doc widht
\usepackage[hyphens]{url}

\usepackage[pdfborder={000}]{hyperref} %referencia

% Ajustar tabla al text width
\usepackage{adjustbox}

% Bibliografia
\usepackage[numbers]{natbib}

% Quitar recuadros rojos del índice
\hypersetup{
    pdfborder = {0 0 0}
}

% Añadir más espacio entre líneas en las tablas
\renewcommand{\arraystretch}{1.5}

% Specific position for images
\usepackage{float}

% Glosario
\usepackage[acronym]{glossaries}

% ********************************************************************
% Re-usable information
% ********************************************************************
\newcommand{\myTitle}{Revisión y mejora de técnicas de hacking para dispositivos NFC\xspace}
\newcommand{\myDegree}{Grado en ...\xspace}
\newcommand{\myName}{José Luis París Reyes\xspace}
\newcommand{\myProf}{Gabriel Maciá Fernández\xspace}
% \newcommand{\myOtherProf}{Nombre Apllido1 Apellido2 (tutor2)\xspace}
%\newcommand{\mySupervisor}{Put name here\xspace}
\newcommand{\myFaculty}{Escuela Técnica Superior de Ingenierías Informática y de
Telecomunicación\xspace}
\newcommand{\myFacultyShort}{E.T.S. de Ingenierías Informática y de
Telecomunicación\xspace}
\newcommand{\myDepartment}{Departamento de PREGUNTAR A MACIÁ\xspace}
\newcommand{\myUni}{\protect{Universidad de Granada}\xspace}
\newcommand{\myLocation}{Granada\xspace}
\newcommand{\myTime}{\today\xspace}
\newcommand{\myVersion}{Version 0.1\xspace}


\hypersetup{
pdfauthor = {\myName joseluis404@correo.ugr.es},
pdftitle = {\myTitle},
pdfsubject = {},
pdfkeywords = {palabra_clave1, palabra_clave2, palabra_clave3, ...},
pdfcreator = {LaTeX con el paquete ....},
pdfproducer = {pdflatex}
}

\makeglossaries

%\hyphenation{}


%\usepackage{doxygen/doxygen}
%\usepackage{pdfpages}
\usepackage{url}
\usepackage{colortbl,longtable}
\usepackage[stable]{footmisc}
%\usepackage{index}

%\makeindex
%\usepackage[style=long, cols=2,border=plain,toc=true,number=none]{glossary}
% \makeglossary

% Definición de comandos que me son tiles:
%\renewcommand{\indexname}{Índice alfabético}
%\renewcommand{\glossaryname}{Glosario}

\pagestyle{fancy}
\fancyhf{}
\fancyhead[LO]{\leftmark}
\fancyhead[RE]{\rightmark}
\fancyhead[RO,LE]{\textbf{\thepage}}
\renewcommand{\chaptermark}[1]{\markboth{\textbf{#1}}{}}
\renewcommand{\sectionmark}[1]{\markright{\textbf{\thesection. #1}}}

\setlength{\headheight}{1.5\headheight}

\newcommand{\HRule}{\rule{\linewidth}{0.5mm}}
%Definimos los tipos teorema, ejemplo y definición podremos usar estos tipos
%simplemente poniendo \begin{teorema} \end{teorema} ...
\newtheorem{teorema}{Teorema}[chapter]
\newtheorem{ejemplo}{Ejemplo}[chapter]
\newtheorem{definicion}{Definición}[chapter]

\definecolor{gray97}{gray}{.97}
\definecolor{gray75}{gray}{.75}
\definecolor{gray45}{gray}{.45}
\definecolor{gray30}{gray}{.94}

\lstset{ frame=Ltb,
     framerule=0.5pt,
     aboveskip=0.5cm,
     framextopmargin=3pt,
     framexbottommargin=3pt,
     framexleftmargin=0.1cm,
     framesep=0pt,
     rulesep=.4pt,
     backgroundcolor=\color{gray97},
     rulesepcolor=\color{black},
     %
     stringstyle=\ttfamily,
     showstringspaces = false,
     basicstyle=\scriptsize\ttfamily,
     commentstyle=\color{gray45},
     keywordstyle=\bfseries,
     %
     numbers=left,
     numbersep=6pt,
     numberstyle=\tiny,
     numberfirstline = false,
     breaklines=true,
   }
 
% minimizar fragmentado de listados
\lstnewenvironment{listing}[1][]
   {\lstset{#1}\pagebreak[0]}{\pagebreak[0]}

\lstdefinestyle{CodigoC}
   {
	basicstyle=\scriptsize,
	frame=single,
	language=C,
	numbers=left
   }
\lstdefinestyle{CodigoC++}
   {
	basicstyle=\small,
	frame=single,
	backgroundcolor=\color{gray30},
	language=C++,
	numbers=left
   }

 
\lstdefinestyle{Consola}
   {
    basicstyle=\scriptsize\bf\ttfamily,
    backgroundcolor=\color{gray30},
    frame=single,
    numbers=none
   }

\lstdefinestyle{Python}
   {
    basicstyle=\small,
	frame=single,
	backgroundcolor=\color{gray30},
	language=Python,
	numbers=left
   }


\newcommand{\bigrule}{\titlerule[0.5mm]}


%Para conseguir que en las páginas en blanco no ponga cabecerass
\makeatletter
\def\clearpage{%
  \ifvmode
    \ifnum \@dbltopnum =\m@ne
      \ifdim \pagetotal <\topskip
        \hbox{}
      \fi
    \fi
  \fi
  \newpage
  \thispagestyle{empty}
  \write\m@ne{}
  \vbox{}
  \penalty -\@Mi
}
\makeatother

\usepackage{pdfpages}
\begin{document}
\begin{titlepage}
 
 
\newlength{\centeroffset}
\setlength{\centeroffset}{-0.5\oddsidemargin}
\addtolength{\centeroffset}{0.5\evensidemargin}
\thispagestyle{empty}

\noindent\hspace*{\centeroffset}\begin{minipage}{\textwidth}

\centering
\includegraphics[width=0.9\textwidth]{imagenes/logo_ugr.png}\\[1.4cm]

\textsc{ \Large TRABAJO FIN DE MÁSTER\\[0.2cm]}
\textsc{ MÁSTER PROPIO EN CIBERSEGURIDAD}\\[1cm]
% Upper part of the page
% 
% Title
{\Huge\bfseries Revisión y mejora de técnicas de hacking para dispositivos NFC\\
}
\noindent\rule[-1ex]{\textwidth}{3pt}\\[3.5ex]
%{\large\bfseries Subtitulo del Proyecto}
\end{minipage}

\vspace{2.5cm}
\noindent\hspace*{\centeroffset}\begin{minipage}{\textwidth}
\centering

\textbf{Autor}\\{José Luis París Reyes}\\[2.5ex]
\textbf{Directores}\\{Gabriel Maciá Fernández}\\[2cm]
\includegraphics[width=0.3\textwidth]{imagenes/etsiit_logo.png}\\[0.1cm]
\textsc{Escuela Técnica Superior de Ingenierías Informática y de Telecomunicación}\\
\textsc{---}\\
Granada, Julio de 2023
\end{minipage}
%\addtolength{\textwidth}{\centeroffset}
%\vspace{\stretch{2}}
\end{titlepage}



\chapter*{}
%\thispagestyle{empty}
%\cleardoublepage

\thispagestyle{empty}

%\begin{titlepage}
 
 
\setlength{\centeroffset}{-0.5\oddsidemargin}
\addtolength{\centeroffset}{0.5\evensidemargin}
\thispagestyle{empty}

\noindent\hspace*{\centeroffset}\begin{minipage}{\textwidth}

\centering
%\includegraphics[width=0.9\textwidth]{imagenes/logo_ugr.jpg}\\[1.4cm]

%\textsc{ \Large PROYECTO FIN DE MÁSTER\\[0.2cm]}
%\textsc{ MÁSTER PROPIO EN CIBERSEGURIDAD DE LA UGR}\\[1cm]
% Upper part of the page
% 

 \vspace{3.3cm}

 \vspace{0.5cm}

%si el proyecto tiene logo poner aquí
\includegraphics{imagenes/logo.png} 
 \vspace{0.5cm}

% Title

{\Huge\bfseries Título del proyecto\\
}
\noindent\rule[-1ex]{\textwidth}{3pt}\\[3.5ex]
%{\large\bfseries Subtítulo del proyecto.\\[4cm]}
\end{minipage}

\vspace{2.5cm}
\noindent\hspace*{\centeroffset}\begin{minipage}{\textwidth}
\centering

\textbf{Autor}\\ {Nombre Apellido1 Apellido2 (alumno)}\\[2.5ex]
\textbf{Directores}\\
{Nombre Apellido1 Apellido2 (tutor1)\\
Nombre Apellido1 Apellido2 (tutor2)}\\[2cm]
%\includegraphics[width=0.15\textwidth]{imagenes/tstc.png}\\[0.1cm]
%\textsc{Departamento de Teoría de la Señal, Telemática y Comunicaciones}\\
%\textsc{---}\\
%Granada, mes de 201
\end{minipage}
%\addtolength{\textwidth}{\centeroffset}
\vspace{\stretch{2}}

 
\end{titlepage}






\cleardoublepage
\thispagestyle{empty}

\begin{center}
{\large\bfseries Revisión y mejora de técnicas de hacking para dispositivos NFC}\\
\end{center}
\begin{center}
José Luis París Reyes\\
\end{center}

%\vspace{0.7cm}
\noindent{\textbf{Palabras clave}: palabra\_clave1, palabra\_clave2, palabra\_clave3, ......}\\

\vspace{0.7cm}
\noindent{\textbf{Resumen}}\\

Poner aquí el resumen.
\cleardoublepage


\thispagestyle{empty}


\begin{center}
{\large\bfseries Project Title}\\
\end{center}
\begin{center}
First name, Family name (student)\\
\end{center}

%\vspace{0.7cm}
\noindent{\textbf{Keywords}: Keyword1, Keyword2, Keyword3, ....}\\

\vspace{0.7cm}
\noindent{\textbf{Abstract}}\\

Write here the abstract in English.

\chapter*{}
\thispagestyle{empty}

\noindent\rule[-1ex]{\textwidth}{2pt}\\[4.5ex]

Yo, \textbf{Nombre Apellido1 Apellido2}, alumno del {\bf Máster Propio en Ciberseguridad} de la {\bf Universidad de Granada}, con DNI XXXXXXXXX, autorizo la ubicación de la siguiente copia de mi {\it Trabajo Fin de Máster} en la biblioteca del centro para que pueda ser consultada por las personas que lo deseen.

\vspace{6cm}

\noindent Fdo: Nombre Apellido1 Apellido2

\vspace{2cm}

\begin{flushright}
Granada a X de mes de 201 .
\end{flushright}


\chapter*{}
\thispagestyle{empty}

\noindent\rule[-1ex]{\textwidth}{2pt}\\[4.5ex]

D. \textbf{Nombre Apellido1 Apellido2 (tutor1)}, Profesor del Departamento YYYY de la Universidad de Granada.

\vspace{0.5cm}

D. \textbf{Nombre Apellido1 Apellido2 (tutor2)}, Profesor del Departamento YYYY de la Universidad de Granada.


\vspace{0.5cm}

\textbf{Informan:}

\vspace{0.5cm}

Que el presente trabajo, titulado \textit{\textbf{Título del proyecto}},
ha sido realizado bajo su supervisión por \textbf{Nombre Apellido1 Apellido2 (alumno)}, y autorizamos la defensa de dicho trabajo ante el tribunal
que corresponda.

\vspace{0.5cm}

Y para que conste, expiden y firman el presente informe en Granada a X de mes de 201 .

\vspace{1cm}

\textbf{Los directores:}

\vspace{5cm}

\noindent \textbf{Nombre Apellido1 Apellido2 (tutor1) \ \ \ \ \ Nombre Apellido1 Apellido2 (tutor2)}

\chapter*{Agradecimientos {\it (opcional)}}
\thispagestyle{empty}

       \vspace{1cm}


Poner aquí agradecimientos...


%\frontmatter

\tableofcontents
\listoffigures
\listoftables
\lstlistoflistings
%
%\mainmatter
%\setlength{\parskip}{5pt}

\chapter{Introducción}

\section {Motivación}

La evolución de la tecnología ha propiciado un gran avance para el ser humano y su calidad de vida, desde los primeros circuitos electrónicos y tarjetas perforadas en los años 60 y 70, hasta el supercomputador más sofisticado en la actualidad, capaz de hacer cálculos de gran complejidad en un tiempo inapreciable, identificación de patrones en ADN o predicción de fenómenos metereológicos, entre otras cosas. \newline

Toda esta tecnología se compone a grandes rasgos de un hardware y un software que le permite funcionar correctamente y satisfacer las necesidades de los usuarios. Dicha evolución ha propiciado la sofisticación de estos dos componentes, provocando que se vuelvan lo suficientemente complejos como para que no se puedan tener en cuenta todas las casuísticas por las que pueden fallar. \newline

El problema se vuelve aún mayor cuando esta tecnología, que permite almacenar, transmitir y procesar información y que está al alcance de cualquiera en forma de un dispositivo portable, es vulnerable, provocando que un usuario con intenciones malintencionadas obtenga información que no debería o adquiera el poder de realizar acciones suplantando la identidad de otra persona. \newline

Además, la popularización de la tecnología \textit{\acrshort{nfc}} que aparece en 2006 ha agravado esta situación, permitiendo establecer una comunicación inalámbrica de corto alcance y alta frecuencia para el intercambio de datos entre dispositivos \cite{NFCFechaCreacion}. Es decir, esta tecnología se encuentra en todas partes, desde los pagos con tarjetas de crédito o débito hasta la sincronización entre dispositivos móviles o la transferencia de información entre ellos. \newline

Es por ello que en este trabajo se revisarán las técnicas actuales que existen contra esa tecnología, se trabajará en posibles mejoras y se hará una realización práctica de las mismas en un laboratorio. La finalidad se basa en mostrar el peligro que supone la vulneración de su seguridad y lo importante que es la continua mejora en la protección de esta tecnología.

\section{Contexto}
%
%\input{capitulos/02_Análisis y Planificacion}
%
%\input{capitulos/03_Diseno}
%
%\input{capitulos/04_Implementacion}
%
%\input{capitulos/05_Evaluación}
%
%\input{capitulos/08_Conclusiones y Trabajo Futuro}
%
%
%%\nocite{*}

%\appendix
%\input{apendices/manual_usuario/manual_usuario}
%%\input{apendices/paper/paper}
%\input{glosario/entradas_glosario}
% \addcontentsline{toc}{chapter}{Glosario}
% \printglossary

\appendix
\newacronym{nfc}{NFC}{Near-Field Communication}
\input{glosario/entradas_glosario}

% Bibliografía
\bibliographystyle{acm}
\bibliography{bibliografia/bibliografia}\addcontentsline{toc}{chapter}{Bibliografía}

\printglossary[type=\acronymtype, title=Acrónimos, toctitle=Acrónimos]\addcontentsline{toc}{chapter}{Acrónimos}

\chapter*{}
\thispagestyle{empty}

\end{document}
